\section{Exact source alternatives on the weak-\texorpdfstring{$B$}{B} II face}
\label{sec:r8-weak-ii-source}
This is an exhaustive source reduction, not a new generic all-tests
admission. Use the original positive source (GE1), with
$z,y,B\in\Inf_{>0}$, $B=H+L$, $0<H,r<1$, $L>0$,
\[
 a_c=1+zr,\quad S=a_c+L,\quad j=a_c-L/y>0,
 \quad G=B+z(H+r-1)>0,
\]
and $B-zyr>0$, $n\geq0$. Write $T=B+zr$ and $g=G/B$.
The already proved common-II source entry gives $\alpha$ weak on
this entire face. All comparisons below are in the pure source field.
The notation $f\asymp h$ means that both $f/h$ and $h/f$ are bounded.

\paragraph{Uniform estimates before any shape specialization.}
The exact positive formula for $\delta$ gives
\[
 \frac\delta{B+z}\sim1.
\]
Moreover $T\geq B$ and $T>z(1-H)$, by $G>0$, while $T\leq B+z$.
Consequently $T/\delta$ is a positive bounded unit, uniformly without
a bound on $z/B$. Since $a_c,S\sim1$, the original formulas imply
\[
 r(2n+1)=\frac{Ta_cS}{\delta}\asymp1,
 \qquad \frac{AH}{y}=\frac{a_cTS}{2\delta}\asymp1.
                                                               \tag{WII1}
\]
Also $0<g\leq1+z$, because $G\leq B+zH$ and $H\leq B$.
The exact identity $z(1-r)=B+zH-G$ rewrites the original $m$ as
\[
 m=\frac{(1+y)a_cS}{2}\frac z\delta
       \frac{1+zH/B-g}{g}.                       \tag{WII2}
\]
If $g$ is a unit, this is bounded. If $g$ is infinitesimal, then
$z(1-r)/B\sim1$, so $z/(B+z)$ is a positive unit and
$m\asymp1/g$. In particular there is no omitted intermediate
large-$m$ case caused by an independent $z/B$ rate.

\paragraph{Exactly which original coefficients are unbounded.}
The following five equivalences hold on the complete weak-$B$ face:
\[
\begin{array}{c|l}
\text{coefficient is unbounded}&\text{necessary and sufficient source condition}\\ \hline
n&r\in\Inf\\
m&G/B\in\Inf\\
A&H/y\in\Inf\\
C&L/(H+y)\in\Inf\\
D&y/B\in\Inf\ \text{or}\ j\in\Inf.
\end{array}                                                     \tag{WII3}
\]
The first three were proved in (WII1)--(WII2). For $C$, use
\[
 M=(1+z)T(H+y)-z(1+y)(1-r)L.
\]
It bounds $C$ above by a bounded multiple of $(H+y)/L$. If
$L/(H+y)$ is infinitesimal, then $z/T$ is bounded and the subtracted
term is infinitesimal relative to $T(H+y)$. Thus in that case
$C\asymp(H+y)/L$ is unbounded. This proves both directions.

For $D$, if $B/y$ is infinitesimal then $L/y$ is infinitesimal and
$j\sim1$. The original formula directly gives $D=O(B/y)$, hence
$D$ weak. If $B/y$ is noninfinitesimal, then $zyr/B$ is infinitesimal,
and the exact positive numerator gives
\[
 D\asymp B/(yj).
\]
Here $j$ is bounded; $j$ infinitesimal forces $L/y\sim1$ and hence
$B/y$ noninfinitesimal. These alternatives give precisely the last
line of (WII3). This proof never replaces $B-zyr$ by $B$ in the
region where that replacement need not be valid.

\paragraph{The positive-total-scale intersection is explicit.}
Put $a=H/B\in(0,1)$ and $\tau=y/B>0$. Every original coefficient is
bounded if and only if all six quantities
\[
 \tau,\quad a,\quad1-a,\quad j,\quad r,\quad g
 \quad\text{are positive bounded units}.          \tag{WII4}
\]
Indeed if $\tau$ is infinitesimal, then $C,D$ are unbounded; if
$\tau$ is unbounded, then $C,A$ are unbounded. With $\tau$ a unit,
(WII3) says respectively that boundedness of $A,C,D,n,m$ is
equivalent to $a,1-a,j,r,g$ being units. Thus (WII4) is exact and
includes arbitrary $z/B$. The original positive-total-scale theorem
already admits this full six-unit interior. After intersecting with
$\delta_{\rm tot}\in\Inf_{>0}$, at least one condition in (WII4)
must fail, equivalently at least one line of (WII3) must occur.

The shape alternatives can also be read as follows:
\[
\begin{array}{c|c|c|c}
 y/B&C&A&D\\ \hline
 \Inf&\text{unbounded}&\text{weak}&\text{unbounded}\\
 \text{unit}&\text{unbounded iff }1-a\in\Inf&
 \text{unbounded iff }a\in\Inf&
 \text{unbounded iff }j\in\Inf\\
 \text{unbounded}&\text{unbounded}&\text{unbounded}&\text{weak}.
\end{array}                                                     \tag{WII5}
\]
In every row $n,m$ are governed independently by $r,g$ as in (WII3).
The exact source inequalities and $n\geq0$ are still imposed; the
table does not assert independent realizability of arbitrary shapes.

\paragraph{The two-weak intersection and the four-product exit.}
There is a stronger exact equivalence:
\[
 \boxed{\ C\in\Inf\quad\Longleftrightarrow\quad
 r,g,H/y,j\in\Inf
 \quad\Longleftrightarrow\quad n,m,A,D\text{ all unbounded}.\ }
                                                               \tag{WII6}
\]
To prove the forward direction, the positive identity
\[
 M=(1+z)TH+yG+yzj+zrL(1+y)
\]
shows $H/L$ infinitesimal. Since $C$ is bounded, (WII3) makes
$L/(H+y)$ a unit. Hence $L\asymp y$, $H/y$ infinitesimal, and
$B\asymp y$. The same four positive summands now imply
$G/\delta,zr/\delta,zj/\delta$ infinitesimal. The quantity
$z/\delta$ cannot be infinitesimal: otherwise $B/\delta\sim1$
and $G/\delta\sim1$, a contradiction. Thus $r,j$ are infinitesimal.
The inequality $G>0$ then gives $B/z$ a positive-unit lower bound;
$z/\delta$ a unit supplies its upper bound. Therefore $g=G/B$ is
infinitesimal too.

Conversely $r,g\in\Inf$ imply, by $G=B+z(H+r-1)$,
\[
 (1-g)B/z=1-H-r,\qquad B/z\sim1.
\]
If also $H/y,j\in\Inf$, then $L/y=a_c-j\sim1$ and
$B/y\sim1$. All four positive summands in $M/(L\delta)$ are now
infinitesimal, so $C$ is weak. Finally (WII3) gives one implication
of the last equivalence. For its converse, $A$ unbounded makes
$H/y$ infinitesimal and $B/y$ bounded, so $D$ unbounded must mean
$j$ infinitesimal. The $n,m$ statements give $r,g\in\Inf$.

The already admitted both-complementary-events-weak theorem removes
the first intersection in (WII6). On the remaining directed face
$C$ strong and $\alpha$ weak, at least one of $n,m,A,D$ is therefore
bounded. Its product with $\alpha$ is infinitesimal. Consequently
the four reciprocal-product exit
\[
 \alpha n,\quad\alpha m,\quad\alpha A,\quad\alpha D
 \quad\text{all strong}
\]
is impossible throughout this directed weak-$\alpha$/strong-$C$
II face. This is a uniform statement, not a failed-product example.

\paragraph{What this does and does not remove.}
The exact all-unit region (WII4) is already admitted by the
positive-total-scale theorem, and the exact four-large intersection
(WII6) is already admitted by the two-weak theorem. Individual real
rational source curves remain covered by the separate curve theorem.
The latter does not automatically admit these generic multiparameter
opening faces. Adjacent dominant-event results require their actual
normalized coefficient support and the stated testing-angle hypotheses;
two unbounded adjacent coefficients alone do not supply those premises.
Thus this reduction identifies the genuine source alternatives still
requiring a full-parent argument. It does not establish that their
intersection with the entire global complement is nonempty or empty.
